\uuid{5NpV}
\titre{Éléments simples}
\auteur{bodin ; hanani ; mzari}
\organisation{exo7}
\licence{CC BY-NC-SA 4.0}
\datecreate{2020-01-20}
\niveau{L1}
\module{Algèbre}
\chapitre{Polynôme, fraction rationnelle}
\sousChapitre{Fraction rationnelle}
\difficulte{2}
\theme{Fractions rationnelles, éléments simples}
\exo7qcmid{lille-1.polynomes.028}

\contenu{
\begin{question}
\qclassification{105.04}
\qtags{motcle=éléments simples}
Soit $P(X)=3X$ et $Q(X) = (X-2)(X-1)^2(X^2-X+1)$.
On écrit
$$\frac{P(X)}{Q(X)} = \frac{a}{X-2} + \frac{b}{X-1} +  \frac{c}{(X-1)^2}
+ \frac{dX+e}{X^2-X+1}.$$
Quelles sont les affirmations vraies ?
\begin{answers}
    \bad{En multipliant par $X-2$, puis en évaluant en $X=2$, j'obtiens $a=1$.}

    \good{En multipliant par $(X-1)^2$, puis en évaluant en $X=1$, j'obtiens $c=-3$.}

    \good{En multipliant par $X$, puis en faisant tendre $X \to +\infty$, j'obtiens la relation $a+b+d=0$.}

    \bad{En évaluant en $X=0$, j'obtiens la relation $a+b+c+e=0$.}
\end{answers}
\begin{explanations}
$\frac{P(X)}{Q(X)} = \frac{3X}{(X-2)(X-1)^2(X^2-X+1)}
=\frac{2}{X-2} + \frac{-3}{X-1} +  \frac{-3}{(X-1)^2}
+ \frac{X+1}{X^2-X+1}$.
\end{explanations}
\end{question}
}
